    \section{Annihilating Polynomials}
        \subsection{Definition of polynomial of a linear operator}
            \hskip \parindent
            \par Let $V$ be a f.d.v.s. over $\mathbb{F}$ and $f$ is a l.op. . Given $p(x) = \sum \limits _{i=0}^m a_i x^i \in \mathbb{F}[x]$, then $P(f) = \sum \limits _{i=0}^m a_i f^i$ where $f^0 = \operatorname{Id}_V$ and $f^i = f \circ f^{i-1}$ for $i \geq 1$ is a linear operator.
            \par Suppose that $p(x)$, $q(x)$ $\in \mathbb{F}[x]$ are s.t. $p(f) = \mathrm{O} = q(f)$. Then $(p+q)(f) = p(f) + q(f) = \mathrm{O}$.
            \par Moreover, if $h(x) \in \mathbb{F}[x]$, then $(h p)(f) = h(f) \cdot p(f) = \mathrm{O}$.
            \par Consider $L_f = \{ p(x) \in \mathbb{F}[x] \mid p(f) = \mathrm{O} \} $
            \begin{itemize}
                \item [(i)] $0 \in L_f$: $0(f)(v) = \mathrm{O}(v) = 0$
                \item [(ii)] $p+q \in L_f$: $p,q\in L_f$, $(p+q)(f) = p(f) + q(f) = \mathrm{O}$
                \item [(iii)] $h p \in L_f$: $p \in L_f$, $(h\cdot p)(f) = h(f) p (f) = \mathrm{O}$
                \item [(iv)] $L_f \neq \mathbb{F}[x]$: Let $p(x) = 1$, then $p(f) = \operatorname{Id}_V$
            \end{itemize}
            \par Annihilating polynomial of $f$ is a non-zero polynomial $p(x) \in \mathbb{F}[x]$ s.t. $p(f) = \mathrm{O}$.
        \subsection{Example of polynomial of l.op.}
            \hskip \parindent
            \par $V$ a f.d.v.s. $p(x) = 2 +6x+7x^2-9x^3$, $q(x) = x-1$
            \par $p(\operatorname*{id}) = 6 \operatorname*{id}$, $q(\operatorname*{id}) = \mathrm{O}$
            \par $q(x) = x -1 \in L_{\operatorname*{id}}$
            \par $r(x) = x \notin L_{\operatorname*{id}}$
        \subsection{Lemma: existence of annihilating polynomial}
            \hskip \parindent
            \par Let $V$ ve a f.d.v.s. over $\mathbb{F}$ and let $f$ be a l.op.. There exists a polynomial $p(x) \in \mathbb{F}[x]$, $p(x) \neq 0$ s.t. $p(f) = \mathrm{O}$.
            \par Proof see [104016 L-251229]-P9.
        \subsection{Remark: minimal degree of annihilating polynomial}
            \hskip \parindent
            \par For $L_f \neq \{0\}$
            \par Let $l = \operatorname*{min} \{\operatorname*{deg}(p(x)) \mid p(x) \in L_f \backslash \{0\}\}$. This minimum exists since $l \leq n^2$.
        \subsection{Proposition: existence and uniqueness of certain degree annihilating polynomial}
            \hskip \parindent
            \par Let $V$ ve a f.d.v.s. over $\mathbb{F}$ and $f$ be a l.op.
            \par There exist a unique polynomial $m(x)$ in $L$ of degree $l$. That is $m(f) = \mathrm{O}$ and $\operatorname*{deg} m = l$.
            \par Proof see [104016 L-251229]-P11
        \subsection{Definition of minimal polynomial}
            \hskip \parindent
            \par Let $V$ be a f.d.v.s. over $\mathbb{F}$ and let $f$ be a l.op. . The unique monic polynomial $m(x)$ s.t.
            \begin{itemize}
                \item [(i)] $m(x)$ is a monic polynomial
                \item [(ii)] $m(f)=\mathrm{O}$
                \item [(iii)] $\operatorname*{deg} m \leq \operatorname*{deg} p(x)$ $\forall p \in \mathbb{F}[x]$ s.t. $p(f)=\mathrm{O}$
            \end{itemize}
            \par is called the minimal polynomial of $f$.
        \subsection{Remark: properties of minimal polynomial}
            \hskip \parindent
            \begin{itemize}
                \item [(1)] In a similar way, the minimal polynomial of a matrix $A \in \mathbb{F}^{n\times n}$ is defined.
                \item [(2)] Similar matrices have the same minimal polynomial. (Proof see [104016 L-251229]-P14)
                \item [(3)] Ex: If $A = [f]_B$ where $B$ is an ordered basis for $V$ and $f$ is a l.op. then $m_f (x) = m_A(x)$.
                \item [(4)] If $p(x) \in \mathbb{F}[x]$, $p(f)=\mathrm{O}$ iff $m(x) \mid p(x)$. (Proof see [104016]-P16)
            \end{itemize}
        \subsection{Example of minimal polynomial}
            \subsubsection{$f= \operatorname*{id}$}
                \hskip \parindent
                \par $f = \operatorname*{id} \Rightarrow m_f (x) = x + a_0$
                \par $\mathrm{O} = m_f (f) = \operatorname*{id} + a_0 \operatorname*{id} = (1 + a_0) \operatorname*{id} \Rightarrow a_0 = -1$
                \par Thus $m_f (x) = x -1$.
            \subsubsection{$A = \begin{bmatrix} 0&-1\\1&0\end{bmatrix}$}
                \hskip \parindent
                \par $A = \begin{bmatrix} 0&-1\\1&0\end{bmatrix}$, if $p(x) = x + a_0$, then $p(A) = A + a_0 \operatorname*{id} = \begin{bmatrix} a_0 & -1 \\ 1 & a_0 \end{bmatrix} \neq \mathrm{O}$
                \par Thus, $\operatorname*{deg}m_A \geq 2$.
                \par $A^2 = \begin{bmatrix} -1 & 0 \\ 0 & -1 \end{bmatrix} = -\operatorname*{id}$
                \par Thus, $A^2 + I = \mathrm{O}$ and $m_A (x) = x^2 + 1$.
        \subsection{Lemma: eigenvalue and minimal polynomial}
            \hskip \parindent
            \par Let $V$ be a f.d.v.s. over $\mathbb{F}$ and let $f$ be a l.op. on $V$. If $\lambda \in \mathbb{F}$, $v \in V$ s.t. $f(v) = \lambda v$, then for all $P \in \mathbb{F}[x]$, $P(f)(v) = P(\lambda) (v)$.
            \par Proof see [104016 L-251231]-P3.
        \subsection{Theorem: minimal polynomial and characteristic polynomial share same roots}
            \hskip \parindent
            \par Let $f : V \rightarrow V$ be a l.op. on an n-dimensional v.s. $V$ over $\mathbb{F}$. $\chi_f(x)$ and $m_f(x)$ have the same root.
            \par Proof see [104016 L-251231]-P4.
        \subsection{Corollary: minimal polynomial of diagonalizable l.op.}
            \hskip \parindent
            \par If $f$ is diagonalizable, with distinct eigenvalues $\lambda_1, \lambda_2, \ldots, \lambda_k$, then (iff) $m_f (x) = (x - \lambda_1)(x - \lambda_2) \cdots (x - \lambda_k)$.
            \par Proof see [104016 L-251231]-P7.
        \subsection{Corollary: inverse of invertible l.op. as polynomial of l.op.}
            \hskip \parindent
            \par Let $V$ be a f.d.v.s. over $\mathbb{F}$ and $f:V \rightarrow V$ be a l.op. 
            \par If $f$ is invertible, then there exists a polynomial $p(x) \in \mathbb{F}[x]$ s.t. $f^{-1} = p(f)$.
            \par Proof see [104016 L-260105]-P13.
            
    \subsection{Cayley-Hamilton Theorem}
        \subsubsection{Statement of Cayley-Hamilton Theorem}            
            \hskip \parindent
            \par Let $V$ be a f.d.v.s. over $\mathbb{F}$ and let $f$ be a l.op. on $V$. Then
                \begin{equation}
                    m_f(x) \mid \chi_f(x) \equiv \chi_f(f) = \mathrm{O}
                \end{equation}
        \subsubsection{Theorem: diagonalizability and minimal polynomial}
            \hskip \parindent
            \par Let $V$ be a f.d.v.s. over $\mathbb{F}$ and let $f$ be a l.op. on $V$. Then $f$ is diagonalizable iff $m_f(x) = \prod\limits_{i=1}^k (x - \lambda_i)$ where $\lambda_1, \lambda_2, \ldots, \lambda_k$ are distinct eigenvalues of $f$.
            \par Proof see [104016 L-251231]-P10.
        \subsubsection{Example of non-diagonalizable matrix via minimal polynomial}
            \hskip \parindent
            \par $A = \begin{bmatrix}1&0&0&0&0\\1&1&0&0&0\\0&0&-1&0&0\\0&0&0&-1&0\\0&0&0&0&3\end{bmatrix}$, $\chi_A (x) = (x-1)^2 (x+1)^2 (x-3)$, $m_A (x) = (x-1)^2 (x+1) (x-3)$
            \par $(A-I)(A+I)(A-3I) \neq \mathrm{O}$
            \par $(A-I)^2(A+I)(A-3I) = \mathrm{O}$
            \par $A$ is not diagonalizable.
    \subsection{Proof for Cayley-Hamilton Theorem}
        \hskip \parindent
        \par Let $f:V \rightarrow V$ be a l.op. on a f.d.v.s. $V$ over $\mathbb{F}$. Then $m_f(x) \mid \chi_f(x) \equiv \chi_f(f) = \mathrm{O}$.
        \subsubsection{Lemma 1: cyclic vector and Cayley-Hamilton Theorem}
            \hskip \parindent
            \par Let $V$ be a n-d.v.s. over $\mathbb{F}$ and let $f$ be a l.op. on $V$. If there exist a vector $v \in V$ s.t. the set $\{ v, f(v), f^2(v), \ldots , f^{n-1}(v) \}$ is a basis of $V$, then $\chi_f(f) = \mathrm{O}$. 
            \par Proof see [104016 L-260105]-P2.
        \subsubsection{Remark: polynomial of block matrix}
            \hskip \parindent
            \par If $A = \begin{bmatrix} A_1&B \\ 0 & A_2\end{bmatrix} \in \mathbb{F}^{n \times n}$ where $A_1 \in \mathbb{F}^{n_1 \times n_1}$, $A_2 \in \mathbb{F}^{n_2 \times n_2}$, and $n=n_1+n_2$. Then for any polynomial $p(x) \in \mathbb{F}[x]$, we have $P(A) = \begin{bmatrix} P(A_1) & C \\ 0 & P(A_2) \end{bmatrix}$ for some matrix $C \in \mathbb{F}^{n_1 \times n_2}$.
            \par Let $p(x) = a_0 + a_1 x + \cdots + a_m x^m \in \mathbb{F}[x]$. Then
            \begin{align*}
                P(A) &= a_0 I_n + a_1 A + a_2 A^2 + \cdots + a_m A^m \\
                &= \begin{bmatrix} a_0 I_{n_1} & 0 \\ 0 & a_0 I_{n_2} \end{bmatrix} + a_1 \begin{bmatrix} A_1 & B \\ 0 & A_2 \end{bmatrix} + a_2 \begin{bmatrix} A_1^2 & A_1 B + B A_2 \\ 0 & A_2^2 \end{bmatrix} + \cdots + a_m \begin{bmatrix} A_1^m & * \\ 0 & A_2^m \end{bmatrix} \\
                &= \begin{bmatrix} P(A_1) & C \\ 0 & P(A_2) \end{bmatrix}
            \end{align*}
        \subsubsection{Lemma 2: block matrix and Cayley-Hamilton Theorem}
            \hskip \parindent
            \par Let $A = \begin{bmatrix} A_1 & B \\ 0 & A_2 \end{bmatrix}$ in $\mathbb{F}^{n \times n}$ where $A_1 \in \mathbb{F}^{n_1 \times n_1}$, $A_2 \in \mathbb{F}^{n_2 \times n_2}$, and $n = n_1 + n_2$. $\chi_{A_1}(A_1) =0$ and $\chi_{A_2}(A_2)=0$, then $\chi_A(A) = 0$.
            \par Proof see [104016 L-260105]-P8.
        \subsubsection{Proof of Cayley-Hamilton Theorem}
            \hskip \parindent
            \par Let $V$ be a n-d.v.s. over $\mathbb{F}$ and let $f$ be a l.op. on $V$. Then $\chi_f(f) = \mathrm{O}$.
            \par Proof see [104016 L-260105]-P9.
